% Part: normal-modal-logic
% Chapter: filtrations
% Section: introduction

\documentclass[../../../include/open-logic-section]{subfiles}

\begin{document}

\olfileid{nml}{fil}{int}

\olsection{प्रस्तावना}

तर्कशास्त्राविषयी एक महत्त्वाचा प्रश्न नेहमीच
असा असतो: ते निर्णेय आहे का? म्हणजे,
“हे !!{formula} वैध आहे का?” या प्रश्नाचे
उत्तर देणारी प्रभावी पद्धत आहे का?
विधानीय तर्कशास्त्र निर्णेय आहे:
सत्यता-कोष्टक रचून एखादे
!!a{formula} सर्वतः-सत्य आहे का हे
आपण प्रभावीपणे तपासू शकतो;
दिलेल्या !!{formula} साठी असे
सत्यता-कोष्टक सांत असते.
पण मोडल सूत्र सर्व प्रतिमानांत सत्य
आहे का हे अशाच रीतीने तपासता
येईल असे स्पष्ट नाही, कारण
प्रतिमाने अनंत संख्येने आहेत.
दिलेल्या सूत्राशी संबंधित सर्व
सांत प्रतिमानांची यादी करता येते;
कारण प्रत्यक्ष त्या !!{formula} मध्ये
येणाऱ्या !!{propositional variable}s
ना जगांचे कोणते उपसंच नेमले
आहेत, एवढेच महत्त्वाचे असते.
प्राप्यता संबंध निश्चित असेल,
तर $V(p)$ ची संभाव्य वेगवेगळी
नेमणूक म्हणजे~$W$ चे सर्व
उपसंच; $\card{W} = n$ असेल,
तर असे $2^n$ उपसंच असतात.
आपल्या !!{formula}~$!A$ मध्ये
$m$ !!{propositional
  variable}s असतील, तर~$n$
जगे असलेली $2^{nm}$ वेगवेगळी
प्रतिमाने मिळतात. प्रत्येक
प्रतिमानात~$!A$ सर्व जगांत
सत्य आहे का हे प्रत्येक जगात
$!A$ चे सत्यतामूल्य काढून
तपासता येते. अर्थात सर्व
संभाव्य प्राप्यता संबंधही
तपासावे लागतात; पण~$n$
जगांवर असे संबंधही सांत
संख्येनेच असतात (नेमके
$W \times W$ चे उपसंच,
म्हणजे $2^{n^2}$ संबंध).

आपल्याला \Log{K} ऐवजी
प्रतिमानांच्या एखाद्या वर्गाने
परिभाषित तर्कशास्त्रात रस
असेल (उदाहरणार्थ,
परावर्ती संक्रमक प्रतिमानांत),
तर प्राप्यता संबंध योग्य प्रकारचा
आहे का हेही तपासता यायला हवे.
आपल्याला हव्या असलेल्या चौकटी
मोडल !!{formula}s ने परिभाष्य
असतील तेव्हा हे करता येते
(उदाहरणार्थ, चौकटीत \Ax{T}
आणि \Ax{4} वैध आहेत का
हे तपासून). म्हणून सर्व सांत
चौकटी क्रमाने घ्याव्यात,
प्रत्येक चौकट इच्छित वर्गात
आहे का हे तपासावे,
त्या चौकटीवरील सर्व संभाव्य
प्रतिमाने नोंदवावीत आणि
प्रत्येकात $!A$ सत्य आहे का
हे तपासावे. नसेल, तर थांबावे:
$!A$ हा संबंधित प्रतिमानांच्या
वर्गात वैध नाही.

या कल्पनेत अडचण आहे:
पाहणे केव्हा थांबवता येईल,
किंवा कधी थांबवता येईलच का,
हे आपल्याला माहीत नाही.
त्या !!{formula} ला सांत
प्रतिप्रतिमान असेल, तर आपली
पद्धत ते शोधेल. पण सांत
प्रतिप्रतिमान नसेल, तर
उत्तर मिळणार नाही.
ते !!{formula} वैध असू शकते
(म्हणजे कोणतेच प्रतिप्रतिमान
नसते), किंवा त्याला केवळ
अनंत प्रतिप्रतिमान असू शकते,
ज्यापर्यंत आपली तपासणी
कधीच पोहोचणार नाही.
प्रतिप्रतिमान असलेल्या प्रत्येक
!!{formula} ला सांत प्रतिप्रतिमानही
असते हे दाखवता आले, तर
ही अडचण दूर होऊ शकते.
असे असेल, तर त्या तर्कशास्त्राला
\emph{सांत प्रतिमान गुणधर्म}
आहे असे म्हणतात.

पण तर्कशास्त्राला सांत प्रतिमान
गुणधर्म आहे हे कसे दाखवायचे?
एक मार्ग म्हणजे~$!A$ चे
अनंत (प्रति)प्रतिमान सांत
प्रतिमानात रूपांतरित करण्याची
पद्धत शोधणे. तसे करता आले,
तर जेव्हा एखाद्या प्रतिमानात
$!A$ सत्य नसते, तेव्हा
त्यातून मिळालेल्या सांत
प्रतिमानातही $!A$ सत्य
नसेल. हे सांत प्रतिमान
सर्व सांत प्रतिमानांच्या
यादीत येईल; आणि वैध
नसलेल्या प्रत्येक सूत्राची
अवैधता आपण अखेरीस ठरवू.
मात्र सूत्र वैध असेल, तर
आपली पद्धत थांबणार नाही.
याशिवाय या पद्धतीतून
मिळणाऱ्या सांत प्रतिमानाच्या
आकाराला कमाल मर्यादा आहे
आणि ती मर्यादा केवळ
!!{formula}~$!A$ वर अवलंबून
आहे हे दाखवता आले, तर
प्रतिप्रतिमान शोधताना
किती आकारापर्यंत सांत
प्रतिमाने तपासायची हे कळेल.
तेवढ्यापर्यंत प्रतिप्रतिमान
सापडले नाही, तर कोणतेच
नाही. अशा वेळी आपली
पद्धत कोणत्याही
!!{formula}~$!A$ साठी
“$!A$ वैध आहे का?” हा
प्रश्न खरोखर निर्णीत करेल.

अनंत रचनांचे सांत रचनांत
रूपांतर करण्याची अनेकदा
यशस्वी होणारी रणनीती म्हणजे
संबंधित बाबतीत सारखे
वागणारे रचनेचे !!{element}s
“एकरूप” मानणे. प्रतिमान~$\mModel{M}$
मध्ये संबंधित बाबतीत सारखे
वागणारी अनंत जगे असतील,
तर ती \emph{सर्व} जगे अशा
जगांच्या सांत संख्येतील
(कदाचित अनंत सदस्य असलेल्या)
“वर्गां”त विभागता येतील.
दुसऱ्या शब्दांत, जगांच्या
संचाचे योग्य प्रकारे विभाजन
करायचे: प्रत्येक वर्गात
कदाचित अनंत जगे असतील,
पण वर्गांची संख्या सांत असेल.
मग नवे प्रतिमान~$\mModel{M^*}$
परिभाषित करू; त्याची जगे
म्हणजे हे वर्ग. जुन्या
प्रतिमानातील सांत संख्येतील
वर्गांमुळे नव्या प्रतिमानात
सांत संख्येने जगे असतील,
म्हणजे सांत प्रतिमान मिळेल.
जग~$w$ ज्या वर्गात आहे
त्याला $[w]$ म्हणू.
जुन्या प्रतिमानात~$!A$ चे
प्रतिप्रतिमान असेल, तर नवे
प्रतिमानही प्रतिप्रतिमान
हवे असल्याने
$\mSat{M}{!A}[w]$ असेल
तेव्हा आणि तेव्हाच
$\mSat{M^*}{!A}[{[w]}]$
असावे. यासाठी विभाजन
आणि~$\mModel{M^*}$ चा
प्राप्यता संबंध योग्य
प्रकारे परिभाषित करावा लागतो.

हे कसे घडेल ते पाहण्यासाठी
प्रथम प्रत्येक जग प्रत्येक
जगाला प्राप्य आहे असे
समजू. मग
$\mSat{M}{\Box !B}[w]$
असेल तेव्हा आणि तेव्हाच
एखाद्या $v \in W$ साठी
$\mSat{M}{\Box !B}[v]$;
$\mModel{M^*}$ साठीही
$[w]$ आणि $[v]$ वापरून
तसेच होईल. पहिली कल्पना
म्हणून दोन जगे~$u$ आणि~$v$
एकाच वर्गात, म्हणजे
समतुल्य, असतील तेव्हा आणि
तेव्हाच प्रतिमान~$\mModel{M}$
मधील सर्व !!{propositional variable}s
वर त्यांचे सत्यत्व एकसारखे
असेल, अशी अट घेऊ.
म्हणजे $\mSat{M}{p}[u]$
असेल तेव्हा आणि तेव्हाच
$\mSat{M}{p}[v]$.
$V^*(p) = \Setabs{[w]}{\mSat{M}{p}[w]}$
ठेवू. आपले ध्येय असे
दाखवणे आहे की
$\mSat{M}{!A}[w]$
असेल तेव्हा आणि तेव्हाच
$\mSat{M^*}{!A}[{[w]}]$.
हे आपण विगमनानेच सिद्ध
करू. आधार प्रकरण
$!A \equiv p$ असेल.
प्रथम $\mSat{M}{p}[w]$
असे समजा. मग व्याख्येनुसार
$[w] \in V^*(p)$;
म्हणून $\mSat{M^*}{p}[{[w]}]$.
आता $\mSat{M^*}{p}[{[w]}]$
असे समजा. याचा अर्थ
$[w] \in V^*(p)$;
म्हणजे~$w$ शी समतुल्य
अशा एखाद्या~$v$ साठी
$\mSat{M}{p}[v]$.
पण “$w$ हे~$v$ शी
समतुल्य आहे” याचा अर्थ
“$w$ आणि~$v$ येथे
नेमकी तीच
!!{propositional variable}s
सत्य आहेत”; म्हणून
$\mSat{M}{p}[w]$.
आता विगमन-पायरी पाहू;
उदाहरणार्थ $!A
\ident \lnot !B$.
मग $\mSat{M}{\lnot !B}[w]$
असेल तेव्हा आणि तेव्हाच
$\mSat/{M}{!B}[w]$;
असेल तेव्हा आणि तेव्हाच
$\mSat/{M^*}{!B}[{[w]}]$
(विगमन गृहीतकानुसार);
असेल तेव्हा आणि तेव्हाच
$\mSat{M^*}{\lnot !B}[{[w]}]$.
इतर गैर-मोडल कारकांसाठीही
तसेच. $\Box$ साठीही हे
चालते: $\mSat{M^*}{\Box
  !B}[{[w]}]$ असे समजा.
याचा अर्थ प्रत्येक $[u]$
साठी $\mSat{M^*}{!B}[{[u]}]$.
विगमन गृहीतकानुसार प्रत्येक
$u$ साठी $\mSat{M}{!B}[u]$.
म्हणून $\mSat{M}{\Box !B}[w]$.

सर्वसाधारण प्रकरणात
$\mModel{M^*}$ साठी
प्राप्यता संबंधही परिभाषित
करावा लागतो; तेव्हा
गोष्टी अधिक गुंतागुंतीच्या
होतात. प्रतिमान~$\mModel{M^*}$
च्या प्राप्यता संबंध~$R^*$
ने वरील विगमन-सिद्धता
चालण्यासाठी आवश्यक अटी
पूर्ण केल्या, तर त्या
प्रतिमानाला
\emph{निस्यंदन} म्हणू.
मग कोणतेही निस्यंदन~$\mModel{M^*}$
येथे $[w]$ वर~$!A$ सत्य
असेल तेव्हा आणि तेव्हाच
$\mModel{M}$ येथे~$w$
वर~$!A$ सत्य असेल.
पण अशी निस्यंदने
\emph{अस्तित्वात आहेत}
हेही दाखवावे लागते;
म्हणजे आवश्यक अटी
पूर्ण करणारा~$R^*$
परिभाषित करता आला पाहिजे.
हे घडण्यासाठी~$u$ आणि~$v$
ही जगे केवळ सर्व
!!{propositional variable}s
वर एकमत असतील तेव्हाच
नव्हे, तर~$!A$ च्या सर्व
उप-!!{formula}s वर
एकमत असतील तेव्हा
समतुल्य मानावी लागतात.
$!A$ ची उप-!!{formula}s
सांत संख्येनेच असल्याने
निस्यंदन सांत राहील.
निस्यंदन परिभाषित
करण्याचा एकच मार्ग नाही;
निस्यंदनाचा प्राप्यता संबंध
आवश्यक गुणधर्म
(उदाहरणार्थ परावर्तित्व,
संक्रमकता इत्यादी)
राखेल यासाठी~$R^*$
ची व्याख्या कल्पकतेने
करावी लागते.


\end{document}
